\chapter{Background and Related Work}
\label{chap:background}

\section{Website Generation Benchmarks}
WebGen-Bench~\cite{lu2025webgenbench} contains 101 website requirements from 20 application domains, each paired with a set of functional test cases (647 in total). A generated project is installed and started, and a \gls{gui} agent based on WebVoyager~\cite{he2024webvoyager} executes each test case and returns \emph{Pass}, \emph{Partial} or \emph{Fail}. A separate \gls{vlm} grades the appearance of the landing page from 1 to 5. The authors also release WebGen-LM, generators fine-tuned on agent trajectories; with the Bolt.diy scaffold, WebGen-LM-7B reaches 28.4\% accuracy and an appearance score of 2.5. Larger general models such as Qwen2.5-Coder-32B score lower (9.5\%), which makes the 7B model a good subject: it is neither at the floor nor at the ceiling of the benchmark.

The benchmark paper reports each configuration as a single run without error bars; its NeurIPS checklist states that error bars were not reported because it would be too computationally expensive~\cite{lu2025webgenbench}. The limitations it lists itself are that generation covers only TypeScript, JavaScript, CSS and HTML, and that only supervised fine-tuning was used to improve open models. Because we compare two arms whose difference may be small relative to run-to-run noise, we repeat runs and report paired confidence intervals (\Cref{sec:order-design}).

WebGen-Bench is a static benchmark. It does not define any interaction between tasks, so neither memory across tasks nor the order of tasks is part of its protocol. Introducing memory makes the order a treatment variable, which we handle explicitly in \Cref{sec:order-design}.

\section{Memory for LLM Agents}
\paragraph{Experience and insight memory.} ExpeL~\cite{zhao2024expel} lets an agent collect trajectories, extract natural-language insights, and maintain them with add, edit, upvote and downvote operations. ReasoningBank~\cite{reasoningbank2025} distils reusable reasoning strategies from both successful and failed trajectories. Generative Agents~\cite{park2023generative} rank memories by recency, importance and relevance.

\paragraph{Skill and code memory.} Voyager~\cite{wang2023voyager} stores executable skills as code with a description and retrieves them by description. Agent Workflow Memory~\cite{wang2025awm} stores reusable workflows. Metis~\cite{dai2026metis} combines textual memory with code tools: plans that recur are promoted to validated, callable tools that the agent executes. Its controlled study with a GPT-4o executor is directly relevant to this work: when memory is built only from earlier tasks (streaming), code memory falls to 53.3\% task completion while text memory reaches 73.3\%, and the authors describe code as efficient but brittle under task variation. Our code memory is of a different kind: fragments placed in the prompt for the generator to imitate, not tools that are executed.

\paragraph{Which form of memory transfers.} Memory Transfer Learning~\cite{kim2026mtl} compares four representations of coding-agent memory, from raw trajectories to abstract insights, with strong models (GPT-5-mini, DeepSeek V3.2, Qwen3-Coder-480B). Abstract insights transfer best across domains, while concrete trajectories can transfer negatively. Neither study tests small models. Our results suggest that the preferred form depends on the capability of the model that has to use the memory: for a 7B generator, abstract textual lessons were rarely executed and became harmful as they accumulated, while short, verified code fragments were followed.

\paragraph{Cross-task reuse on WebGen-Bench.} Super Library Agent~\cite{sung2026superlibrary} also generates WebGen-Bench applications in sequence, maintaining a shared library of reusable components. Its goal is maintainability without loss of functionality: on 24 tasks in three suites drawn from single topical clusters, with a strong backbone, accuracy changes from 76.04\% to 77.21\%, a difference that is not statistically significant, while code size and token use fall. It does not store lessons from failures and does not vary task order; the authors note that single-topic suites are a favourable setting for reuse. We study memory of failures and successes for correctness, with a small generator, on all 101 tasks under the official protocol, and treat task order as a variable.

\paragraph{Memory maintenance.} Mem0~\cite{chhikara2025mem0} decides at write time whether a new memory should be added, should update an old one, or is redundant. Memento~\cite{zhou2025memento} learns the utility of stored cases from their outcomes. Our hygiene component borrows these ideas: deduplication at storage time and a utility score for each memory.

\paragraph{Task order in continual agent evaluation.} Benchmarks for agents that learn across tasks report that the order and composition of the task stream affect results. Evo-Memory~\cite{wei2025evomemory} finds that baselines vary substantially between easy-to-hard and hard-to-easy sequences, and AgentCL~\cite{shu2026agentcl} shows that memory methods can only be distinguished on streams that expose opportunities for reuse. We therefore compare orders in which tasks of the same domain are adjacent with orders in which they are mixed (\Cref{sec:order-design}).

\section{Small Models, Long Prompts and Instruction Following}
Several results suggest that adding text to a prompt is not free. Models become less likely to satisfy all instructions as their number grows~\cite{harada2025curse,jaroslawicz2025ifscale}; information in the middle of a long context is used less than at its start or end~\cite{liu2024lost}; irrelevant context distracts models~\cite{shi2023distracted}; and longer inputs alone can reduce performance even with perfect retrieval~\cite{du2025contextlength}. Models can be trained to prioritise instructions by source~\cite{wallace2024hierarchy}, which motivates keeping the user requirement explicitly above any memory content. Small models also struggle to benefit from long reasoning chains~\cite{li2025smallmodels}, which is consistent with our negative result for a short planning step (\Cref{sec:history}).

\section{Self-Correction and Debugging}
Reflection on failed attempts is related to self-correction methods such as Self-Refine~\cite{madaan2023selfrefine}, Self-Debug~\cite{chen2023selfdebug} and Chain-of-Verification~\cite{dhuliawala2023cove}. These methods repair the \emph{same} task with feedback. We do not allow this in the formal comparison, because repairing a benchmark task using its hidden evaluation results would leak the answer. Reflection is used only to write memory for \emph{later} tasks. The investigation steps we ask the reflection model to follow are modelled on systematic debugging~\cite{zeller2009whyprograms}.
